% ===================================================================
%  図表第 15 号 — 市場。— 食べもの。
%
%  Ceci n'est PAS une transcription : c'est une traduction, faite en
%  2026, en japonais d'aujourd'hui. Voir texto/ja/10-zuhyo-01.tex
%  pour la consigne, qui vaut pour les seize fichiers.
%
%  LES NOMS DE FROMAGE ET DE LIQUEUR RESTENT CEUX DE L'ORIGINAL, en
%  katakana : オランダ, カマンベール, グリュイエール, キルシュ,
%  キュラソー. Le japonais nomme ainsi ces produits-la, et le
%  vocabulaire du livret est justement celui des choses etrangeres
%  qu'un marche apporte.
% ===================================================================

%%K t15-apar-1 apar
\VUsaut{4.0mm}
\VUtitre{15.0pt}{60}{図表第 15 号}
\VUsaut{3.0mm}

%%K t15-tit-1 sub
\VUpk{11.4pt}{40}{市場。--- 食べもの。}
\VUsaut{3.0mm}

%%K t15-01-1 p
1. --- 私たちの食物にいる品を届けてくれる商人の大半は、\VUgras{屋根つき市場}の
\VUgras{大屋根} \textsuperscript{(1)} の下か、近くの街路に集まっている。

%%K t15-02-1 p
2. --- \VUgras{豚肉屋} \textsuperscript{(2)} は\VUgras{ハム} \textsuperscript{(3)}、
\VUgras{血の腸詰} \textsuperscript{(4)}、\VUgras{ソーセージ} \textsuperscript{(5)}、
\VUgras{もつの腸詰} \textsuperscript{(6)}、\VUgras{塩豚} \textsuperscript{(7)} を売って
いて、\VUgras{バターとチーズを売る女} \textsuperscript{(8)} のとなりに立っている。その
店先には赤い皮の\VUgras{オランダチーズ} \textsuperscript{(9)}、\VUgras{カマンベール}
\textsuperscript{(10)}、大きな車輪の形の\VUgras{グリュイエール} \textsuperscript{(11)}、
新しい\VUgras{バターの塊} \textsuperscript{(12)} が並んでいる。

%%K t15-03-1 p
3. --- その向かいでは\VUgras{八百屋の女} \textsuperscript{(13)} が客に立派な
\VUgras{トマト} \textsuperscript{(14)}、\VUgras{カリフラワー} \textsuperscript{(15)}、
\VUgras{サラダ菜} \textsuperscript{(16)}、\VUgras{キャベツ} \textsuperscript{(17)}、
\VUgras{かぼちゃ} \textsuperscript{(19)}、\VUgras{メロン} \textsuperscript{(18)}、
\VUgras{きのこ} \textsuperscript{(20)} まですすめている。ところが\VUgras{ビール屋の若い者}
が\VUgras{配達車} \textsuperscript{(21)} を——\VUgras{ビール樽} \textsuperscript{(22)} を
山と積んだ車を——ちょうどその店先に停めたので、客が近づけない。百姓の女が
\VUgras{アーティチョーク} \textsuperscript{(23)} を一籠、彼女に届けにきた。

%%K t15-tit-2 sub
\VUsaut{4.0mm}
\VUpk{10.2pt}{40}{売り買い。}
\VUsaut{3.0mm}

%%K t15-04-1 p
4. --- 市場は活気にあふれている。\VUgras{競り} \textsuperscript{(24)} の時刻がきたので
ある。買い物にくる\VUgras{主婦} \textsuperscript{(26)} たちは、\VUgras{もつ屋の女}
\textsuperscript{(25)} の店先を通りかかると足をとめ、\VUgras{胃袋のもつ}か
\VUgras{子牛の頭}を買う。

%%K t15-05-1 p
5. --- 太った\VUgras{料理人} \textsuperscript{(27)} が\VUgras{魚屋の女}
\textsuperscript{(28)} を相手に、たいそう立派な\VUgras{まぐろ}の値を掛け合っている。
その女は気の向くままに売り値を吊り上げる。彼女のところには\VUgras{うなぎ、したびらめ、
ます}、\VUgras{こい}が大量に入っている。\VUgras{たら}も\VUgras{むらさきいがい}も
たいそう新しい。

%%K t15-tit-3 sub
\VUsaut{4.0mm}
\VUpk{10.2pt}{40}{安い品と高い品。}
\VUsaut{3.0mm}

%%K t15-06-1 p
6. --- そのそばに店を出した\VUgras{鶏肉屋の女} \textsuperscript{(29)} はたいそう
律義である。\VUgras{七面鳥}、\VUgras{がちょう}、\VUgras{あひる}、ただの
\VUgras{にわとり}を売るとき、正しい値しか取らない。

%%K t15-07-1 p
7. --- 広場の角の二階には、少し前から\VUgras{布屋} \textsuperscript{(30)} が店を
かまえている。そこでは買い手はいつでも、たいそう立派な\VUgras{食卓布}、
\VUgras{ナプキン、前掛け}、\VUgras{ふきん}のそろいを見つけられる。同じ階には
\VUgras{刃物屋} \textsuperscript{(31)} が仕事場を開いた。市場の女商人たちの\VUgras{包丁}を
研ぎ、百姓のためにいちばん硬い\VUgras{鋼}で\VUgras{剪定鎌}をこしらえる。

%%K t15-tit-4 sub
\VUsaut{4.0mm}
\VUpk{10.2pt}{40}{食料品店。}
\VUsaut{3.0mm}

%%K t15-08-1 p
8. --- その仕事場の下には、少し前から大きな食料品店が開いている。もはや昔の、
質素で目立たない\VUgras{乾物屋} \textsuperscript{(32)} ——日用の品、すなわち\VUgras{茶}
\textsuperscript{(33)}、\VUgras{チョコレート} \textsuperscript{(34)}、\VUgras{角砂糖}
\textsuperscript{(37)}、\VUgras{石鹸} \textsuperscript{(38)} しか売らなかった店——では
ない。そこではどんな品も売っている。\VUgras{堅焼きの菓子}、\VUgras{缶詰}
\textsuperscript{(35)}、\VUgras{リキュール} \textsuperscript{(36)}、\VUgras{ラム、
コニャック、キルシュ、アニス酒}、\VUgras{キュラソー}。狩りの獲物も——
\VUgras{野うさぎ} \textsuperscript{(39)}、\VUgras{つぐみ} \textsuperscript{(40)}、
\VUgras{やまうずら} \textsuperscript{(41)}、\VUgras{うずら} \textsuperscript{(42)}、
\VUgras{野がも} \textsuperscript{(43)}、\VUgras{きじ} \textsuperscript{(44)} も——
\VUgras{バナナ} \textsuperscript{(46)} や\VUgras{パイナップル}のような南の果物と同じく
たやすく手に入る。

%%K t15-09-1 p
9. --- たいそう愛想のよい店の\VUgras{若い者} \textsuperscript{(45)} が、望みの品を
いつでも出してくれる。すぐりやまるめろの\VUgras{ジャム} \textsuperscript{(47)}、
\VUgras{ラード} \textsuperscript{(48)}、\VUgras{きゅうりの酢漬} \textsuperscript{(49)}、
塩漬の\VUgras{いわし} \textsuperscript{(50)}。

%%K t15-09-2 p
これは\VUgras{女の客} \textsuperscript{(51)} にたいそう便利である。いちばんありふれた品の
すぐそばに、いちばん珍しい初物といちばん味のよい果物——汁の多い\VUgras{梨}
\textsuperscript{(52)}、\VUgras{すもも} \textsuperscript{(53)}、\VUgras{ぶどう}
\textsuperscript{(54)}、\VUgras{桃} \textsuperscript{(55)}、\VUgras{アーモンド}
\textsuperscript{(56)}、\VUgras{くるみ} \textsuperscript{(57)} ——が見つかるのだから。

%%K t15-tit-5 sub
\VUsaut{4.0mm}
\VUpk{10.2pt}{40}{客たち。}
\VUsaut{3.0mm}

%%K t15-10-1 p
10. --- \VUgras{米} \textsuperscript{(58)} と\VUgras{レンズ豆} \textsuperscript{(59)}
は薫製の\VUgras{にしん} \textsuperscript{(60)} の箱のそばにある。\VUgras{じゃがいも}
\textsuperscript{(61)} と\VUgras{いんげん} \textsuperscript{(63)} は\VUgras{りんご}
\textsuperscript{(62)}、\VUgras{いちじく} \textsuperscript{(64)}、\VUgras{干しすもも}
\textsuperscript{(65)}、\VUgras{はしばみの実} \textsuperscript{(66)} と隣り合っている。
この店には新しい\VUgras{卵} \textsuperscript{(67)} も\VUgras{オリーブ}
\textsuperscript{(68)} もあるので、\VUgras{籠} \textsuperscript{(70)} や
\VUgras{買い物網} \textsuperscript{(73)} はみるみるいっぱいになる。

%%K t15-11-1 p
11. --- この立派な店の\VUgras{コーヒー} \textsuperscript{(72)} は\VUgras{女中}
\textsuperscript{(69)} や料理女のあいだで評判が高く、それも当然である。年とった
\VUgras{食料品屋} \textsuperscript{(71)} が自分でこの上なく念入りに炒っているのだから。

%%K t15-tit-6 sub
\VUsaut{4.0mm}
\VUpk{10.2pt}{40}{見ものの数々。}
\VUsaut{3.0mm}

%%K t15-12-1 p
12. --- 誰もが買い物のために市場へ来るわけではない。大屋根へ通じる小道の絵のような
人通りの中には、実にさまざまの人が見える。人のよい\VUgras{屠場の職人}
\textsuperscript{(75)} は\VUgras{手押し車} \textsuperscript{(74)} を前に押している。その
そばには休暇の兵隊が二人、きらびやかな軍服を得意げに見せている。一人は
\VUgras{軽騎兵} \textsuperscript{(76)} で、青い\VUgras{ドルマン}を着て\VUgras{羽根飾りの
筒帽}をかぶり、もう一人は\VUgras{ズアーヴ兵の伍長}で、ふくらんだ袴をはき、頭には
\VUgras{フェズ}をのせている。

%%K t15-13-1 p
13. --- \VUgras{パン屋} \textsuperscript{(80)} は\VUgras{パン焼き店}
\textsuperscript{(78)} の戸口に立っている。\VUgras{パン} \textsuperscript{(79)} の生地を
こねたばかりで、飾り窓に並ぶ\VUgras{クロワッサン} \textsuperscript{(81)}、
\VUgras{小型パン} \textsuperscript{(82)}、\VUgras{丸パン} \textsuperscript{(83)} を窯から
出したところである。

%%K t15-14-1 p
14. --- パン焼き店の上には腕のよい\VUgras{肌着仕立ての女} \textsuperscript{(84)} が
住んでいて、\VUgras{刺繍女} \textsuperscript{(85)} を何人か使っている。そのとなりには
\VUgras{洗濯とアイロンの女} \textsuperscript{(86)} が住んでいる。\VUgras{洗濯物}
\textsuperscript{(88)} を洗って干したあと糊をつけ、\VUgras{アイロン}
\textsuperscript{(87)} でのしをかける。

%%K t15-tit-7 sub
\VUsaut{4.0mm}
\VUpk{10.2pt}{40}{肉、魚、野菜。}
\VUsaut{3.0mm}

%%K t15-15-1 p
15. --- 一階にある\VUgras{肉屋} \textsuperscript{(89)} では、あらゆる\VUgras{肉}を
売っている。\VUgras{羊肉} \textsuperscript{(90)}、\VUgras{牛肉} \textsuperscript{(94)}、
\VUgras{子牛肉} \textsuperscript{(92)} を、\VUgras{羊のもも、ビフテキ、焼き肉}、
\VUgras{あばら肉}にして売るのである。\VUgras{肉屋の若い者} \textsuperscript{(93)} が
これらの肉をみな切り分け、\VUgras{髄の骨} \textsuperscript{(96)} を別にしておく。肉屋の
戸口には\VUgras{牡蠣売りの女} \textsuperscript{(97)} がいて、\VUgras{牡蠣}
\textsuperscript{(98)} を売っている。望めば殻を開けてくれる。

%%K t15-16-1 p
16. --- これらの商人はみな鑑札料か場所代を払っている。だから\VUgras{巡査}
\textsuperscript{(100)} は情け容赦なく、気の毒な\VUgras{野菜の呼び売り女}
\textsuperscript{(99)} を追いまわす。この女たちは手押し車で\VUgras{にんじん、はつかだいこん、
かぶ、にら葱}や\VUgras{アスパラガスの束、新しいえんどう、ほうれん草、すいば、
クレソン}、\VUgras{パセリ}を売り歩き、商人たちの商売敵になっているのである。

%%K t15-tit-8 sub
\VUsaut{4.0mm}
\VUpk{10.2pt}{40}{巡査に気をつけろ！}
\VUsaut{3.0mm}

%%K t15-17-1 p
17. --- \VUgras{にんにく} \textsuperscript{(101)} と\VUgras{玉ねぎ}を売る小僧も、
自分もまた市場のそばで品を売ってはいけないことをよく知っている。彼は逃げ出し、
走りながらもう少しで\VUgras{かに、ざりがに、小えび}の籠をひっくり返すところだった。
\VUgras{いせえび}と\VUgras{オマール}の\VUgras{売り子} \textsuperscript{(102)} の籠である。

%%K t15-18-1 p
18. --- 誰もが動きまわり、たいへんな騒ぎである。それでも、\VUgras{ガラス屋}
\textsuperscript{(103)} が呼び歩く声も、\VUgras{炭屋} \textsuperscript{(104)} の
掛け声もはっきり聞こえる。炭屋は\VUgras{炭の袋} \textsuperscript{(105)} を届けるため、
ちょっとのあいだ荷車を離れたところである。
\VUsaut{6.0mm}
\VUfilet{22mm}
